\documentclass{report}
\usepackage{amsmath,amssymb}
\setlength{\parindent}{0mm}
\setlength{\parskip}{1em}
\begin{document}
\begin{center}
\rule{6in}{1pt} \
{\large Lars Ruthotto \\
{\bf jInv - A Flexible Framework for Parallel PDE Parameter Estimation in Julia}}

Emory University \\ Department of Mathematics and Computer Science \\ 400 Dowman Drive \\ W408 \\ 30322 Atlanta \\ GA \\ USA
\\
{\tt lruthotto@emory.edu}\\
Eldad Haber\end{center}

This talk presents jInv, a new framework for solving parameter estimation
problems involving partial differential equations (PDEs). Parameter
estimation problems of this kind arise frequently in many scientific
disciplines, for example, in medical imaging, geophysical imaging, non
destructive testing and economics. Parameter estimation is an inverse
problem and typically ill-posed. The problem is typically formulated as
an optimization problem with constraints that are given by the PDEs. The
unknowns are parameters of the PDEs, which correspond to physical
properties of the object to be measured. The objective is to minimize the
misfit between PDE simulations and measured data plus some regularization
term.

jInv is implemented in the relatively new programming language
Julia~\cite{BezansonEtAl2012}, which offers a few advantages. First,
Julia is a dynamic language therefore allows for rapid prototyping, is
easy to read, change, distribute, and document. Second, Julia comes with
a just-in-time compiler and therefore the runtime of the code can in many
cases approach that of C or Fortran. Third, Julia provides tools that
simplify parallel and distributed computing. Our code was initially
developed on a standard laptop but supports also homogeneous clusters, as
well as heterogeneous cloud computing services.

In this talk, we consider forward problems of the form
\begin{equation}\label{eq:fwd}
d_{ijk} = p_k^T u_{ij} + e_{ijk},\quad {\rm where}\quad A(m,\omega_i) u_{ij} = q_j,
\end{equation}
for all $\quad \forall i=1,\ldots,n_f, j=1,\ldots,n_s, k=1,\ldots,n_r$.
Here, $A(m,\omega_i)$ is a discretization of the governing PDE for a
discrete model $m$ and additional parameter $\omega_i$, for example,
frequency. Further, $p_k$ represent receivers, $q_j$ the applied sources,
and $u_{ij}$ is the field caused by the $i$-th source and the $j$-th
frequency. The noise, $e_{ijk}$ is assumed to identically distributed
white noise. It is important to note that computing the data requires the
computation of $n_f \cdot n_s$ fields associated with $n_f$ different
linear systems and $n_s$ right hand sides.

In practical applications, the number of sources and frequencies is
generally large, and so a large number of PDE solves is required in the
inversion. Thus, efficient methods for solving the PDEs are essential. In
jInv, one can chose between state-of-the-art iterative or direct methods
depending on the problem size and computational resources. Also, jInv
supports both regular and stretched tensor meshes that can be adapted to
the source locations.

For given data, sources, receivers, and frequencies, the goal of the
parameter estimation problem is to estimate the model $m$. Therefore, we
consider the bound-constrained discrete optimization problem
\begin{equation}\label{eq:inv}
\min_{m} \sum_{ijk} \phi( p_{i}^\top {A}(\omega_{j},m)^{-1}q_{k}
,d_{ijk}) + \alpha R(m)\quad
{\rm subject\ to } \quad m_{\rm low} \leq m \leq m_{\rm up}.
\end{equation}
Here, $\phi$ measures the misfit of the predicted and measured data, $R$
is a regularizer, and $m_{\rm low}$ and $m_{\rm up}$ are lower- and upper
bounds on the components of $m$, respectively. The parameter $\alpha>0$
balances between the misfit and regularizer. In its current version jInv
provides some commonly used misfit functions (weighted $L_2$ and $L_1$
norms) and regularization (diffusion, total variation) that can be
exchanged modularly.

The optimization problem is solved using a projected Gauss-Newton method
that uses a Steepest Descent step to identify the active set. The Hessian
approximation of the data term is implemented in a matrix-free fashion,
since it involves $A(m,\omega_i)^{-1}$ for $i=1,\ldots,n_f$, which are in
general large and dense. Intermediate results of the PDE solver such as a
factorizations or preconditioner are re-used to accelerate computations
with the Hessian. A projected PCG method is used for determining the
search direction and projected Armijo condition is used for line search.

A particular focus of the talk is Julia's potential for parallel and
distributed computing that is used in jInv to reflect the parallel
structure of the misfit term in~\eqref{eq:inv}. Note that all terms can
be computed independently and in parallel. Therefore a variety of
parallelization strategies are supported by jInv that can be tested to
exploit the given computational resources.

We use a distributed data model that divides the data, sources, receivers
and frequencies among all available workers and therefore parallelizes
evaluation of the misfit. One benefit is the reduced memory consumption,
in particular regarding the fields and intermediate results from the PDE
solves that are held in memory by the respective worker. Another benefit
of this model is the reduction of communication costs, since only the
model $m$ as well as vectors of equivalent size must be communicated.

The talk will present numerical examples geophysical imaging. Particular
emphasis is on the parallel implementation using Julia. We will show the
scalability of the code on a single-user laptop, a large parallel
workstation, and a cloud computing engine.

\begin{thebibliography}{1}
\bibitem{haberBook2014}
E.~Haber.
\newblock {\em Computational Methods in Geophysical Electromagnetics}.
\newblock SIAM, Philadelphia, 2014.


\bibitem{BezansonEtAl2012}
Jeff Bezanson, Stefan Karpinski, Viral~B Shah, and Alan Edelman.
\newblock {Julia: A Fast Dynamic Language for Technical Computing}.
\newblock {\em arXiv.org}, 2012.
\end{thebibliography}


\end{document}
